% Part: first-order-logic
% Chapter: natural-deduction
% Section: proof-theoretic-notions

\documentclass[../../../include/open-logic-section]{subfiles}

\begin{document}

\iftag{FOL}
      {\olfileid{fol}{ntd}{ptn}}
      {\olfileid{pl}{ntd}{ptn}}

\olsection{Begrippen over formele afleidingen}

\begin{editorial}
In deze paragraaf verzamelen we de definities van de relatie
``is afleidbaar uit'' en van consistentie voor natuurlijke deductie.
\end{editorial}

\begin{explain}
We hebben al belangrijke begrippen over betekenis gedefinieerd:
geldigheid, logisch gevolg en de vraag of iets waar te maken is. Nu
definiëren we de bijbehorende \emph{bewijstheoretische begrippen},
dus begrippen over formele afleidingen. Daarbij vragen we niet of
!!{sentence}s waar zijn in !!{structure}s. We vragen of bepaalde
!!{sentence}s wel of niet !!{derivable} zijn uit andere. De ontdekking
dat de twee soorten begrippen samenvallen, was belangrijk. De
\emph{correctheidsstelling} en de \emph{volledigheidsstelling} zeggen
dat zij
inderdaad samenvallen.
\end{explain}


\begin{defn}[Stellingen]
!!^a{sentence}~$!A$ is een \emph{stelling} als er in de natuurlijke
deductie !!a{derivation} van~$!A$ bestaat waarin alle aannamen zijn
!!{discharged}. We schrijven $\Proves !A$ als $!A$ een stelling is.
Als dat niet zo is, schrijven we $\Proves/ !A$.
\end{defn}

\begin{defn}[!!^{derivability}]
!!^a{sentence} $!A$ is \emph{!!{derivable} uit} een verzameling
!!{sentence}s~$\Gamma$, genoteerd als $\Gamma \Proves !A$, als er
!!a{derivation} met conclusie~$!A$ bestaat. Elke aanname in die
!!{derivation} moet ofwel !!{discharged} zijn, ofwel in~$\Gamma$
staan. Als $!A$ niet !!{derivable} is uit $\Gamma$, schrijven we
$\Gamma \Proves/ !A$.
\end{defn}

\begin{defn}[Consistentie]
Een verzameling !!{sentence}s~$\Gamma$ is \emph{inconsistent} precies
als $\Gamma \Proves \lfalse$. Is $\Gamma$ niet inconsistent, dus als
$\Gamma \Proves/ \lfalse$, dan noemen we haar \emph{consistent}.
\end{defn}

\begin{prop}[Reflexiviteit]
\ollabel{prop:reflexivity}
Als $!A \in \Gamma$, dan $\Gamma \Proves !A$.
\end{prop}

\begin{proof}
De aanname $!A$ op zichzelf is al !!a{derivation} van~$!A$. De enige
!!{undischarged} aanname is $!A$, en die staat in~$\Gamma$.
\end{proof}
  

\begin{prop}[Monotonie]
\ollabel{prop:monotonicity}
Als $\Gamma \subseteq \Delta$ en $\Gamma \Proves !A$, dan $\Delta
\Proves !A$.
\end{prop}

\begin{proof}
Elke !!{derivation} van $!A$ uit $\Gamma$ is ook !!a{derivation} van
$!A$ uit~$\Delta$.
\end{proof}

\begin{prop}[Transitiviteit]
\ollabel{prop:transitivity}
Als $\Gamma \Proves !A$ en $\{!A\} \cup \Delta \Proves
!B$, dan $\Gamma \cup \Delta \Proves !B$.
\end{prop}

\begin{proof}
Uit $\Gamma \Proves !A$ volgt dat er !!a{derivation}~$\delta_0$
van~$!A$ bestaat. Alle !!{undischarged} aannamen daarvan staan
in~$\Gamma$. Uit $\{!A\} \cup \Delta \Proves !B$ volgt dat er
!!a{derivation}~$\delta_1$ van~$!B$ bestaat. Alle !!{undischarged}
aannamen daarvan staan in~$\{!A\} \cup \Delta$. Beschouw nu:
\begin{prooftree}
  \AxiomC{$\Delta, \Discharge{!A}{1}$}
  \RightLabel{$\delta_1$}
  \DeduceC{$!B$}
  \DischargeRule{\Intro{\lif}}{1}
  \UnaryInfC{$!A \lif !B$}
  \AxiomC{$\Gamma$}
  \RightLabel{$\delta_0$}
  \DeduceC{$!A$}
  \RightLabel{\Elim{\lif}}
  \BinaryInfC{$!B$}
\end{prooftree}
Alle !!{undischarged} aannamen staan nu in $\Gamma \cup
\Delta$. Dit toont dus $\Gamma \cup \Delta \Proves !B$ aan.
\end{proof}

Als $\Gamma = \{!A_1, !A_2, \ldots, !A_k\}$ een eindige verzameling
is, mogen we de korte vorm
$!A_1,!A_2,\ldots,!A_k \Proves !B$ gebruiken voor
$\Gamma \Proves !B$. In het bijzonder betekent
$!A \Proves !B$ dat $\{!A\} \Proves !B$.

Merk op: als $\Gamma \Proves !A$ en $!A
\Proves !B$, dan $\Gamma \Proves !B$. Daaruit volgt ook het volgende.
Als $!A_1, \dots, !A_n \Proves !B$ en $\Gamma \Proves !A_i$ voor
elke~$i$, dan $\Gamma \Proves !B$.

\begin{prop}
\ollabel{prop:incons}
De volgende uitspraken zijn equivalent.
    \begin{enumerate}
      \item \( \Gamma \) is inconsistent.
      \item \( \Gamma \Proves {!A} \) geldt voor elke
        !!{sentence}~\( {!A} \).
      \item Zowel \( \Gamma \Proves {!A} \) als
        \( \Gamma \Proves \lnot {!A} \) geldt voor minstens één
        !!{sentence}~\( {!A} \).
    \end{enumerate}
\end{prop}

\begin{proof}
Oefening.
\end{proof}

\tagprob{FOL}
\begin{prob}
Bewijs \olref[fol][ntd][ptn]{prop:incons}
\end{prob}
\tagendprob

\tagprob{notFOL}
\begin{prob}
Bewijs \olref[pl][ntd][ptn]{prop:incons}
\end{prob}
\tagendprob

\begin{prop}[Compactheid]
  \ollabel{prop:proves-compact}
  \begin{enumerate}
  \item Als $\Gamma \Proves !A$, dan is er een eindige
    deelverzameling $\Gamma_0 \subseteq \Gamma$ waarvoor
    $\Gamma_0 \Proves !A$.
  \item Als elke eindige deelverzameling van~$\Gamma$ consistent is,
    dan is $\Gamma$ consistent.
  \end{enumerate}
\end{prop}

\begin{proof}
  \begin{enumerate}
    \item Als $\Gamma \Proves !A$, is er !!a{derivation}~$\delta$
      van~$!A$ uit~$\Gamma$. Laat $\Gamma_0$ bestaan uit de
      !!{undischarged} aannamen van~$\delta$. Elke !!{derivation} is
      eindig. Daarom kan $\Gamma_0$ maar eindig veel !!{sentence}s
      bevatten. Dus is $\delta$ !!a{derivation} van~$!A$ uit een
      eindige~$\Gamma_0 \subseteq \Gamma$.
    \item Dit is de contrapositie van (1) in het bijzondere geval
      $!A \ident \lfalse$.
  \end{enumerate}
\end{proof}
\end{document}
